\begin{proposition}[Euler characteristic of a normal crossing divisor]\hypertarget{tag-02XY}{}\label{night:P4.6}\label{stmt:P4.6}
\StatementRecord{night:P4.6}{proved}{02XY}{the definitions and constructions stated in the proof}

Let $X$ be a compact complex threefold with $b_2(X)=0$, and let $F$ be a
reduced compact normal-crossings divisor. We allow different local
branches of one global irreducible component to meet. If $T(F)$ is the
number of ordinary triple points, then
\[
 e(F)=T(F),\qquad T(F)\text{ is even}.
\]
For a fibre $\sum_j m_jF_j$ with this reduced support, the divisor
in the formula is $F=\sum_jF_j$.
\end{proposition}
\begin{proof}

\begin{enumerate}[label={$\langle1\rangle$\arabic*.},leftmargin=*]
\item Write $\nu:S=\coprod S_i\to X$ for the normalization immersion, and $D_i$
for the conductor divisor on $S_i$. The surfaces $S_i$ are smooth and
the curves $D_i$ are reduced nodal divisors. The quotient
$N_i=\nu_i^*T_X/T_{S_i}$ is a line bundle. A local equation of $F$ gives
\[
 N_i\simeq \nu_i^*\OO_X(F)\otimes\OO_{S_i}(-D_i).
\]
Indeed, for the coordinate equation $x_1\cdots x_k=0$ the normal
derivative on $x_1=0$ is $x_2\cdots x_k$, whose zero divisor is $D_i$.
These are morphisms of holomorphic vector bundles, including when
$S_i$ is not embedded globally.

\item Since $b_2(X)=0$, both $c_1(T_X)$ and $c_1(\OO_X(F))$ vanish in real
cohomology. Duality gives $b_4(X)=0$, so $c_2(T_X)$ also vanishes. The
tangent exact sequence yields
\[
 K_{S_i}\equiv-D_i,\qquad c_2(S_i)=K_{S_i}^2.
\]
If $a=\sum_i\chi(\OO_{S_i})$, the surface Riemann--Roch formula
$12\chi(\OO_{S_i})=K_{S_i}^2+e(S_i)$ consequently gives
$\sum_i e(S_i)=6a$.

\item Let $\Gamma$ be the normalization of the double-curve locus, with one
sheet for each unordered pair of local branches. The normalization of
$\coprod D_i$ is an unramified degree-two cover of $\Gamma$, recording
the choice of branch. Each triple point contributes three nodes to
$\coprod D_i$. Adjunction gives
$\chi(\OO_{D_i})=-D_i\cdot(D_i+K_{S_i})/2=0$. The normalization
sequence for these nodal curves therefore gives
\[
 0=2\chi(\OO_\Gamma)-3T(F)=e(\Gamma)-3T(F).
\]
The normalization complex of $F$ is the exact alternating restriction
complex of the local coordinate hyperplanes. Permuting the branches
globally replaces its double-curve term by $\OO_\Gamma$ tensored with
the sign line of the two-sheeted branch cover. That line has order at
most two, and hence degree zero on every component; it has the same
coherent Euler characteristic as $\OO_\Gamma$. The triple term has
length one at each triple point. It follows that
\[
 \chi(\OO_F)=a-\chi(\OO_\Gamma)+T(F)=a-T(F)/2.
\]
The exact sequence for the Cartier divisor $F$ and analytic
Hirzebruch--Riemann--Roch give $\chi(\OO_F)=0$, since
$\operatorname{ch}(\OO_X(-F))=\operatorname{ch}(\OO_X)$ in rational
cohomology. Thus $T(F)=2a$. Finally ordinary Euler additivity gives
\[
 e(F)=\sum_i e(S_i)-e(\Gamma)+T(F)=6a-2T(F)=T(F).
\]
At a point with $r=1,2,3$ local branches, the alternating
contribution is $r-\binom r2+\binom r3=1$. This gives the
displayed Euler additivity formula and proves the identity and parity.

\end{enumerate}
\end{proof}